% Source: https://artofproblemsolving.com/community/c3h1592878_bizarre_logic_question
% Stable identity: aops-topic-1592878-q01
\textbf{With at least one person.} ``Not all people are aliens'' means: there is at least one person who is not an alien. Take any person $p$ (one exists by assumption). Since no person is an alien, $p$ is not an alien. So $p$ is a person who is not an alien, and the statement ``Not all people are aliens'' is true.

\textbf{Without that assumption.} The conclusion is no longer guaranteed. In logic, ``all people are aliens'' and ``no person is an alien'' are both statements about every person, and if there are no people at all, both are \emph{vacuously} true: there is no person who could be a counterexample. In that case ``Not all people are aliens'' is false, while ``No person is an alien'' is true.

So: $\boxed{\text{yes, provided at least one person exists}}$; the existence of a person is exactly what is needed.
